
\subsection*{2.6}
\phantomsection
\label{sga3:VIIB:2.6}
\addcontentsline{toc}{subsection}{\SGARefLink{sga3:VIIB:2.6}{2.6}. Lie algebras of formal groups}

We shall now study infinitesimal formal groups, to which the sections
that follow are devoted. In this study, Lie algebras play a fundamental
role.

Suppose first that the base ring \(k\) is Artinian, and let \(G\) be a
formal group over \(k\). The Lie algebra of \(G\) admits three different
interpretations, all of which we shall use:

\medskip
\noindent\textbf{a)}
Let
\[
  D=k[d]/(d^2)
\]
be the algebra of dual numbers over \(k\), and let
\[
  \delta:D\longrightarrow k
\]
be the homomorphism that annihilates \(d\). For every formal group
\(G\) over \(k\), \(\operatorname{Lie}(G)\) is the kernel of
\(G(\delta)\). Thus, by definition, there is an exact sequence of
groups
\begin{center}
  \phantomsection
  \label{sga3:VIIB:2.6:dual-numbers-exact-sequence}
  \(
1\longrightarrow\operatorname{Lie}(G)
\longrightarrow G(D)
\xrightarrow{\ G(\delta)\ }G(k)
\longrightarrow1.
\)

\end{center}

\medskip
\noindent\textbf{b)}
Let \(A\) be the affine algebra of \(G\), and let
\[
  I_A=\ker\varepsilon_A
\]
be its augmentation ideal. The elements of \(G(D)\) are the morphisms
of profinite \(k\)-algebras
\[
  f:A\longrightarrow D.
\]
The condition \(G(\delta)(f)=1\) is equivalent to
\[
  \delta\circ f=\varepsilon_A.
\]
Since
\[
  x-\varepsilon_A(x)\cdot1_A\in I_A
  \qquad\text{for every \(x\in A\),}
\]
this is equivalent to \(f(I_A)\subset k\cdot d\). Hence
\(\ker(f)\) contains \(I_A^2\), and therefore also its closure
\(\overline{I_A^2}\). Thus \(f\) induces a continuous linear map
\[
  f':
  I_A/\overline{I_A^2}
  =
  \omega_{G/k}
  \longrightarrow k
\]
such that, for every \(x\in A\),
\[
  f(x)
  =
  \varepsilon_A(x)\cdot1_D+f'(\bar x)\cdot d,
\]
where \(\bar x\) denotes the image of
\(x-\varepsilon_A(x)\cdot1_A\) in
\(I_A/\overline{I_A^2}\). It follows that the map
\[
  f\longmapsto f'
\]
defines a bijection from \(\operatorname{Lie}(G)\) onto the
topological dual
\[
  \omega_{G/k}^{\dagger}
\]
of \(\omega_{G/k}\) \textup{(cf.\ \SGARefLink{sga3:VIIB:0.2.2}{0.2.2})}.

This bijection respects the group structures. Indeed, let \(f\) and
\(g\) be two elements of \(\operatorname{Lie}(G)\). Their product
\(f\cdot g\) is the composite \(h\circ\Delta_A\), where
\[
  h:A\widehat\otimes A\longrightarrow D
\]
is defined by
\[
  h(b\widehat\otimes b')=f(b)g(b').
\]
If \(a\in I_A\), then, by \SGARefLink{sga3:VIIA:2.1}{2.1}\textup{\SGARefLink{sga3:vii:paragraph:5-1:equation:ii}{(ii)}},
\[
  \Delta_A(a)-a\widehat\otimes1-1\widehat\otimes a
  \in
  I_A\widehat\otimes I_A.
\]
Consequently,
\[
  (f\cdot g)(a)=f(a)+g(a)
\]
\textup{(cf.\ also II, \SGARefLink{sga3:ii:lemma:3-10}{3.10})}.

Henceforth we identify
\[
  \operatorname{Lie}(G)=\omega_{G/k}^{\dagger}
\]
by means of the bijection \(f\mapsto f'\) just described. The group
\(\operatorname{Lie}(G)\) is thereby endowed with a \(k\)-module
structure.

\medskip
\noindent\SGAInlineRefTarget{sga3:VIIB:2.6:item:c}\textbf{c)}
Let \(A^\dagger\) and \(D^\dagger\) be the \(k\)-modules dual to
\(A\) and \(D\), and let
\[
  \{1_D^\dagger,d^\dagger\}
\]
be the basis dual to the basis \(\{1_D,d\}\) of \(D\) over \(k\);
thus \(1_D^\dagger=\delta\).\footnote{Translator's note: the French text prints
\(1_D^*\) and \(d^*\) in the two occurrences immediately below,
after defining \(1_D^\dagger\) and \(d^\dagger\) as the dual basis.
The consistent \(\dagger\)-notation is used here.}
Since \(D\) is free of finite rank over \(k\), the canonical map
\[
  \operatorname{Hom}_c(A,D)
  \longrightarrow
  \operatorname{Hom}_k(D^\dagger,A^\dagger),
  \qquad
  f\longmapsto{}^t f,
\]
is bijective. On the other hand, \({}^t f\) is determined by the
values
\[
  {}^t f(1_D^\dagger)
  \qquad\text{and}\qquad
  {}^t f(d^\dagger)=x.
\]
The condition \(G(\delta)(f)=1\) is equivalent to
\[
  {}^t f(1_D^\dagger)=\varepsilon_A.
\]
Moreover, one sees readily that \(f\) is compatible with
multiplication if and only if \textup{(cf.\ \SGARefLink{sga3:VIIB:2.3}{2.3})}
\[
  \tag{*}
  \label{sga3:VIIB:2.6:primitive-relation}
  \delta_G(x)=\phi_G(x\otimes1+1\otimes x).
\]

Finally, it is clear that a continuous linear map
\[
  f:A\longrightarrow D
\]
that is compatible with multiplication and satisfies
\(\delta\circ f=\varepsilon_A\) sends the unit element of \(A\) to
that of \(D\).
The map \(f\mapsto x\) therefore allows us to identify
\(\operatorname{Lie}(G)\) with the set of ``primitive elements'' of
\(\mathbf H(G)\), that is, the elements \(x\in\mathbf H(G)\)
satisfying relation \textup{(*)}. If \(x\) and \(y\) are two such
elements, then
\[
\begin{aligned}
  \delta_G(xy)
  &=
  \delta_G(x)\delta_G(y)
  \\
  &=
  \phi_G\bigl(
    (x\otimes1+1\otimes x)
    (y\otimes1+1\otimes y)
  \bigr)
  \\
  &=
  \phi_G\bigl(
    xy\otimes1+x\otimes y+y\otimes x+1\otimes xy
  \bigr),
\end{aligned}
\]
whence
\[
  \delta_G(xy-yx)
  =
  \phi_G\bigl(
    (xy-yx)\otimes1+1\otimes(xy-yx)
  \bigr).
\]
This shows that the \(k\)-module \(\operatorname{Lie}(G)\) is
identified with a Lie subalgebra of \(\mathbf H(G)\); we call
\(\operatorname{Lie}(G)\) the \emph{Lie algebra} of
\(G\).\footnote{\label{sga3:VIIB:nde:108}\textit{Editors' note.} Comparing with
VII\textsubscript{A}, \SGARefLink{sga3:VIIB:2.5:proposition}{2.5}, we see that if \(K\) is a group scheme of
finite type over \(k\) and \(G\) is the formal completion of \(K\) at
the identity---that is, if \(\mathcal A(G)\) is the completion of the
local ring \(\mathcal O_{K,e}\) for the \(\mathfrak m\)-adic topology,
where
\[
  \mathfrak m=\ker\bigl(
    \varepsilon:\mathcal O_{K,e}\longrightarrow k
  \bigr)
\]
---then \(\mathbf H(G)\) is identified with the algebra of
distributions \(U(K)\), and \(\operatorname{Lie}(G)\) with
\(\operatorname{Lie}(K)\). The hypothesis that \(K\) is of finite
type over \(k\) is used to ensure that
\(\mathfrak m/\mathfrak m^2\) is a \(k\)-module of finite length, and
hence discrete, so that its topological dual coincides with its
ordinary dual. In particular, when \(G\) is a finite formal
\(k\)-group---that is, when \(\mathcal A(G)\) is a finite
\(k\)-module---in which case \(G\) may also be regarded as the
\(k\)-group scheme \(\operatorname{Spec}\mathcal A(G)\), the two
definitions of \(\operatorname{Lie}(G)\) coincide.}

\subsubsection*{2.6.1}
\phantomsection
\label{sga3:VIIB:2.6.1}

Let \(k\) now be an arbitrary pseudocompact ring, and let \(G\) be a
formal group over \(k\). We call the functor \(\operatorname{Lie}(G)\)
that assigns to every object \(C\) of \(\operatorname{Alf}/k\) the
\(C\)-Lie algebra of the \(C\)-formal group
\[
  G'=G\widehat\otimes_k C
\]
the \emph{\(\mathbf O_k\)-Lie algebra of \(G\)}.\footnote{%
\label{sga3:VIIB:nde:109}The editors expanded the original in what
follows.}
Put
\[
  A'=A\widehat\otimes_k C.
\]
Since \(I_A\) is a direct factor of \(A\), we have
\[
  I_{A'}
  =
  I_A\widehat\otimes_k C
  =
  I_A\widehat\otimes_A A'.
\]
Moreover,
\[
  \omega_{G/k}=I_A\widehat\otimes_A k,
  \qquad
  \omega_{G'/C}=I_{A'}\widehat\otimes_{A'}C.
\]
It follows that
\[
  \omega_{G'/C}
  =
  \omega_{G/k}\widehat\otimes_A C,
\]
and hence
\[
\begin{aligned}
  \operatorname{Lie}(G)(C)
  &=
  \operatorname{Hom}_c\bigl(
    \omega_{G/k}\widehat\otimes_A C,C
  \bigr),
  \\
  \text{i.e.}\qquad
  \operatorname{Lie}(G)
  &=
  \mathbf V_k^f(\omega_{G/k}).
\end{aligned}
\]
With the notation of \SGARefLink{sga3:VIIB:1.2.3.B}{1.2.3.B}, \SGARefLink{sga3:VIIB:1.2.3.E}{Proposition 1.2.3.E} therefore shows
that \(\operatorname{Lie}(G)\) is flat over \(\mathbf O_k\) if and
only if \(\omega_{G/k}\) is a projective pseudocompact \(k\)-module.
